\documentclass[a4paper,12pt]{article}
\usepackage[T2A]{fontenc}
\usepackage[utf8]{inputenc}
\usepackage[russian]{babel}
\usepackage{amsmath,amssymb,mathtools}
\usepackage{helvet}
\renewcommand{\familydefault}{\sfdefault}
\usepackage{icomma}
\usepackage{geometry}
\geometry{left=19mm,right=19mm,top=20mm,bottom=20mm}
\usepackage{microtype}
\usepackage{tikz}
\usetikzlibrary{calc,arrows.meta,positioning,shapes.geometric}
\usepackage{pgfplots}
\pgfplotsset{compat=1.18}
\usepackage{tabularx,booktabs,longtable}
\usepackage{enumitem}
\usepackage{hyperref}
\usepackage{parskip}
\usepackage{fancyhdr}
\usepackage{setspace}
\usepackage{xcolor}
\definecolor{accent}{HTML}{14777A}
\definecolor{darkaccent}{HTML}{164E51}
\definecolor{textgray}{HTML}{353A3A}
\definecolor{muted}{HTML}{737D7E}
\definecolor{gridcolor}{HTML}{D8E0DE}
\hypersetup{hidelinks,pdfauthor={Артём Александрович Лёвин},pdftitle={Чек-лист: показательные функции по графику}}
\setlength{\parindent}{0pt}
\setlength{\parskip}{6pt}
\setlist[itemize]{leftmargin=6mm,itemsep=2pt,topsep=4pt}
\setlist[enumerate]{leftmargin=8mm,itemsep=3pt,topsep=4pt}
\renewcommand{\arraystretch}{1.22}
\setstretch{1.13}
\pagestyle{fancy}
\fancyhf{}
\renewcommand{\headrulewidth}{0pt}
\fancyfoot[L]{\footnotesize\textcolor{muted}{А.\,А. Лёвин \quad / \quad Анна Баннова}}
\fancyfoot[R]{\footnotesize\textcolor{muted}{\thepage}}
\newcommand{\checkitem}{\(\square\)\enspace}
\newcommand{\strong}[1]{\textcolor{darkaccent}{\textbf{#1}}}
\newcommand{\sectionline}{\par\vspace{-3pt}\noindent\textcolor{accent}{\rule{\linewidth}{0.5pt}}\par\vspace{2pt}}
\newcommand{\lessonfig}[1]{\begin{center}\footnotesize\textcolor{muted}{#1}\end{center}}
\begin{document}
\thispagestyle{empty}
\begin{center}
\vspace*{22mm}
{\large\textcolor{accent}{\textbf{МАТЕМАТИКА \ / \ ЕГЭ}}}\par\vspace{20mm}
{\fontsize{31}{35}\selectfont\bfseries ЧЕК-ЛИСТ}\par\vspace{10mm}
{\fontsize{21}{28}\selectfont\bfseries Показательные функции\\ по графику}\par\vspace{9mm}
{\large Метод опорных (узловых) точек}\par
\vspace{16mm}
\textcolor{accent}{\rule{0.58\linewidth}{0.7pt}}\par\vspace{15mm}
{\large\bfseries Анна Баннова}\par\vspace{4mm}
23.10.26\par\vspace{10mm}
\normalsize\textcolor{textgray}{\textbf{Лёвин Артём Александрович}\\[2mm] эксклюзивно для Анны Банновой}
\end{center}
\vfill
\begin{center}
\begin{tikzpicture}
\draw[accent,opacity=.65,line width=1.2pt]
 (0,0) .. controls (3,1.25) and (5,-1) .. (7.8,.2)
       .. controls (10,1.0) and (12,-.6) .. (15.8,.35);
\end{tikzpicture}
\end{center}
\clearpage

\section*{01. Что нужно уметь}
\sectionline
\checkitem Понимать запись $f(t)$: подставить $x=t$ в формулу функции.\par
\checkitem Отличать задачу \(\boldsymbol{f(t)=?}\) от задачи \(\boldsymbol{f(x)=k,\ x=?}\).\par
\checkitem Находить на графике \strong{опорные точки} с удобными координатами.\par
\checkitem Записывать уравнение по каждой точке $(x_0;y_0)$: \(y_0=f(x_0)\).\par
\checkitem Составлять систему для неизвестных $a$ и $b$ и решать её подстановкой.\par
\checkitem Использовать \(a^0=1\) и \((a^m)^n=a^{mn}\).\par
\checkitem Учитывать \(a>0\) и \(a\ne1\) для рассматриваемой показательной функции.\par
\checkitem Сохранять минусы в координатах, показателях и при переносе слагаемых.\par

\subsection*{Опорные точки и запись координат}
\strong{Опорной (узловой) точкой} называем удобную для вычислений точку графика, обычно с целыми координатами. В паре \((x;y)\) \textbf{первым} указывается $x$, \textbf{вторым} --- $y=f(x)$.

Для функций с двумя неизвестными параметрами $a,b$ выбираем \strong{две точки с различными абсциссами}, задающие два уравнения. Особенно удобны точки с $x=0$, $x=1$, $x=-1$ или с $y=1$ (когда вся функция является степенью).

\subsection*{Степени: краткая памятка}
\begin{align*}
&a^0=1\quad(a>0),\qquad a^{-n}=\frac{1}{a^n},\\
&(a^m)^n=a^{mn},\qquad \sqrt[n]{a}=a^{1/n}\quad(a>0),\\
&a^{2k}=c\ (c>0)\ \Longrightarrow\ a=c^{1/(2k)}\quad\text{при }a>0.
\end{align*}
\textbf{Ограничение:} в функции $a^x$ при действительных $x$ берём $a>0$; в стандартной показательной функции также $a\ne1$. Если при решении $a^2=4$, алгебраически получаем $a=\pm2$, но основанием может быть только $a=2$.

\subsection*{Два принципиально разных вида функции}
\begin{center}
\begin{tabularx}{\linewidth}{@{}X X@{}}
\toprule
\(f(x)=a^x+b\) & \(f(x)=a^{x+b}\)\\
\midrule
$b$ складывается \emph{со всей степенью}. & $b$ находится \emph{в показателе}.\\
$x=0$ даёт $f(0)=1+b$. & $f(x)=1$ даёт $x+b=0$ при $a>0$, $a\ne1$.\\
\bottomrule
\end{tabularx}
\end{center}

\clearpage
\section*{02. Пример: \(f(x)=a^x+b\)}
\sectionline
На графике отмечены $A(0;-2)$ и $B(1;-1)$. Найдите $f(6)$.
\begin{center}
\begin{tikzpicture}
\begin{axis}[
 width=12.0cm,height=7.0cm,axis lines=middle,xmin=-2,xmax=2.4,
 ymin=-4,ymax=3.5,xtick={-2,-1,0,1,2},ytick={-3,-2,-1,0,1,2,3},
 grid=both,grid style={gridcolor,very thin},minor tick num=0,
 xlabel={$x$},ylabel={$y$},xlabel style={at={(axis description cs:1,.42)},anchor=west},
 ylabel style={at={(axis description cs:.46,1)},anchor=south},
 tick label style={font=\footnotesize},clip=false]
\addplot[accent,very thick,domain=-2:2.15,samples=200]{2^x-3};
\addplot[only marks,mark=*,mark size=2.7pt,accent] coordinates {(0,-2) (1,-1)};
\node[below left,font=\small] at (axis cs:0,-2) {$A(0;-2)$};
\node[above left,font=\small] at (axis cs:1,-1) {$B(1;-1)$};
\end{axis}
\end{tikzpicture}
\end{center}
\textbf{Шаг 1.} Подставляем координаты двух точек:
\[
\begin{cases}
-2=a^0+b,\\
-1=a^1+b.
\end{cases}
\quad\Longleftrightarrow\quad
\begin{cases}
-2=1+b,\\
-1=a+b.
\end{cases}
\]
\textbf{Шаг 2.} Находим коэффициенты:
\[
 b=-3,\qquad a-3=-1\;\Longrightarrow\;a=2.
\]
\textbf{Шаг 3.} Восстанавливаем формулу и вычисляем:
\[
 f(x)=2^x-3,\qquad f(6)=2^6-3=64-3=\boxed{61}.
\]
\textcolor{muted}{Самопроверка: $f(0)=1-3=-2$ и $f(1)=2-3=-1$ --- обе точки принадлежат найденному графику.}

\clearpage
\section*{03. Та же функция, другой вопрос}
\sectionline
На графике $f(x)=a^x+b$ отмечены точки $A(0;-1)$ и $B(2;1)$. Найдите $x$, если $f(x)=25$.
\[
\begin{cases}
-1=a^0+b,\\ 1=a^2+b
\end{cases}
\ \Longrightarrow\ b=-2,\quad a^2=3\quad\Longrightarrow\quad a=\sqrt3.
\]
\[
 f(x)=(\sqrt3)^x-2=3^{x/2}-2.
\]
Теперь используем условие $f(x)=25$:
\[
 3^{x/2}-2=25\quad\Longleftrightarrow\quad 3^{x/2}=27=3^3
 \quad\Longleftrightarrow\quad \frac{x}{2}=3
 \quad\Longleftrightarrow\quad \boxed{x=6}.
\]
\textbf{Разница вопросов:} запись $f(6)$ означает \emph{известен $x$ --- найдите значение функции}; уравнение $f(x)=25$ означает \emph{известно значение функции --- найдите $x$}. Поиск $a,b$ остаётся тем же.

\subsection*{Когда на графике есть точка с отрицательным $x$}
Для точек $A(0;-3)$ и $B(-2;0)$ получаем
\[
\begin{cases}
-3=a^0+b,\\0=a^{-2}+b
\end{cases}
\Longrightarrow b=-4,\quad a^{-2}=4.
\]
Следовательно, $a^2=\frac14$, а с учётом $a>0$ имеем $a=\frac12=2^{-1}$. Поэтому
\[
 f(x)=2^{-x}-4,\qquad f(-4)=2^4-4=12.
\]
\textbf{Обратите внимание:} при переходе от $(2^{-1})^x$ к $2^{-x}$ минус \emph{не исчезает}.

\subsection*{Частые ошибки}
\begin{itemize}
\item Перепутали $x$ и $y$ у точки; правильно $\left(x;y\right)$ и $y=f(x)$.
\item Потеряли минус у ординаты $-1$ или абсциссы $-2$.
\item Приняли $a^2=4$ за $a=4$; на самом деле $a=2$ при $a>0$.
\item Записали $a^{x}+b$ как $a^{x+b}$: показатели устроены по-разному.
\end{itemize}

\clearpage
\section*{04. Алгоритм: метод опорных точек}
\sectionline
\vspace{6mm}
\begin{center}
\begin{tikzpicture}[
>=Latex,
start/.style={ellipse,draw=darkaccent,thick,fill=accent!6,minimum width=5.5cm,minimum height=9mm,align=center,font=\small},
proc/.style={rectangle,rounded corners=2pt,draw=accent,thick,minimum width=9.6cm,text width=9.0cm,minimum height=11mm,align=center,font=\small},
choice/.style={diamond,aspect=2.45,draw=darkaccent,thick,fill=accent!6,text width=4.0cm,inner sep=1pt,align=center,font=\small},
branch/.style={rectangle,rounded corners=2pt,draw=accent,thick,minimum width=5.5cm,text width=4.9cm,minimum height=16mm,align=center,font=\small},
arrow/.style={-Latex,thick,accent},node distance=10mm]
\node[start] (start) {Дана формула с неизвестными $a,b$ и график};
\node[proc,below=of start] (pick) {Выберите две удобные опорные точки $A(x_1;y_1)$ и $B(x_2;y_2)$, $x_1\ne x_2$};
\node[proc,below=of pick] (sys) {Составьте систему $y_1=f(x_1)$ и $y_2=f(x_2)$};
\node[proc,below=of sys] (simpl) {Упростите степени: $a^0=1$, $a^1=a$; найдите $b$ и $a$, оставив $a>0$};
\node[proc,below=of simpl] (func) {Подставьте $a,b$ в исходную формулу функции};
\node[choice,below=9mm of func] (q) {Что требуется найти?};
\node[branch,below left=18mm and -2.5mm of q] (val) {Значение $f(t)$?\\Подставьте $x=t$ и вычислите};
\node[branch,below right=18mm and -2.5mm of q] (arg) {Аргумент $x$?\\Решите уравнение $f(x)=k$};
\node[start,below=18mm of q,yshift=-34mm] (end) {Проверьте вычисления и запишите ответ};
\draw[arrow] (start)--(pick);
\draw[arrow] (pick)--(sys);
\draw[arrow] (sys)--(simpl);
\draw[arrow] (simpl)--(func);
\draw[arrow] (func)--(q);
\draw[arrow] (q.south west) -| node[pos=.32,above left,font=\footnotesize]{\textbf{значение}} (val.north);
\draw[arrow] (q.south east) -| node[pos=.32,above right,font=\footnotesize]{\textbf{аргумент}} (arg.north);
\draw[arrow] (val.south) |- ([yshift=4mm]end.west) -- (end.west);
\draw[arrow] (arg.south) |- ([yshift=4mm]end.east) -- (end.east);
\end{tikzpicture}
\end{center}
\clearpage

\section*{05. Вторая конструкция: \(f(x)=a^{x+b}\)}
\sectionline
На графике отмечены $A(-3;1)$ и $B(1;4)$. Найдите $f(-7)$.
\begin{center}
\begin{tikzpicture}
\begin{axis}[
 width=11.7cm,height=6.1cm,axis lines=middle,xmin=-4.6,xmax=2,
 ymin=0,ymax=5.9,xtick={-4,-3,-2,-1,0,1,2},ytick={0,1,2,3,4,5},
 grid=both,grid style={gridcolor,very thin},minor tick num=0,
 xlabel={$x$},ylabel={$y$},tick label style={font=\footnotesize},clip=false]
\addplot[accent,very thick,domain=-4.6:2,samples=200]{pow(2,(x+3)/2)};
\addplot[only marks,mark=*,mark size=2.7pt,accent] coordinates {(-3,1) (1,4)};
\node[above left,font=\small] at (axis cs:-3,1) {$A(-3;1)$};
\node[above left,font=\small] at (axis cs:1,4) {$B(1;4)$};
\end{axis}
\end{tikzpicture}
\end{center}
\textbf{Шаг 1.} Подставляем точки именно в $a^{x+b}$:
\[
\begin{cases}
1=a^{-3+b},\\
4=a^{1+b}.
\end{cases}
\]
\textbf{Шаг 2.} Поскольку $1=a^0$ и $a\ne1$, из первого уравнения $-3+b=0$, значит $b=3$.
\[
4=a^{1+3}=a^4\quad\Longrightarrow\quad a=4^{1/4}=2^{1/2}=\sqrt2.
\]
\textbf{Шаг 3.} Восстанавливаем функцию и вычисляем:
\[
 f(x)=(2^{1/2})^{x+3}=2^{(x+3)/2},
\]
\[
 f(-7)=2^{(-7+3)/2}=2^{-2}=\frac14=\boxed{0{,}25}.
\]
\textcolor{muted}{Проверка по точкам: $f(-3)=2^0=1$, $f(1)=2^2=4$.}

\clearpage
\section*{06. Самопроверка перед тренировкой}
\sectionline
\begin{enumerate}[label=\arabic*.]
\item Вы выписали координаты в порядке $(x;y)$ и подставили $y=f(x)$?
\item В показателе стоит $x$ или $x+b$? Проверьте скобки.
\item Вы заменили $a^0$ единицей и правильно сохранили знаки?
\item Если получили $a^2=c$ или $a^4=c$, выбрали положительный корень?
\item При наличии отрицательной степени использовали $a^{-n}=1/a^n$?
\item Вы вернулись к первоначальному вопросу: $f(t)$ или $x$?
\item Подстановка найденных коэффициентов воспроизводит координаты обеих точек?
\end{enumerate}
\vspace{8mm}
\begin{center}
{\Large\bfseries\textcolor{darkaccent}{От опорных точек к ответу}}\\[3mm]
$\text{график}\ \longrightarrow\ \text{система}\ \longrightarrow\ (a,b)
\ \longrightarrow\ \text{формула}\ \longrightarrow\ \text{ответ}$
\end{center}
\clearpage
\section*{07. Тренировочный блок --- Путь А}
\sectionline
\textbf{2 задания.} На каждом графике отмечены опорные точки. Для каждой задачи найдите $a,b$, восстановите формулу функции и ответьте на поставленный вопрос.
\par\vspace{4mm}
\noindent\begin{minipage}[t]{.485\linewidth}
\textbf{Задание 1.}\quad$ f(x)=a^x+b $\\[1mm]
\begin{tikzpicture}
\begin{axis}[
 width=\linewidth,height=5.2cm,scale only axis=false,
 xmin=-1.2,xmax=1.4,ymin=-2.5,ymax=5,
 xtick={-1,0,1},ytick={-2,0,2,4},axis lines=middle,
 grid=both,grid style={gridcolor,very thin},
 tick label style={font=\scriptsize},
 xlabel={$x$},ylabel={$y$},
 label style={font=\scriptsize},
 xlabel style={at={(axis description cs:1.01,.19)},anchor=west},
 ylabel style={at={(axis description cs:.40,1)},anchor=south},
 clip=true]
\addplot[accent,very thick,domain=-1.2:1.4,samples=200]{4^x-2};
\addplot[only marks,mark=*,mark size=2.4pt,accent] coordinates {(0,-1) (1,2)};
\end{axis}
\end{tikzpicture}\\[-1mm]
\footnotesize Точки $A(0;-1)$ и $B(1;2)$.\\
Найдите $f(3)$.
\end{minipage}\hfill
\begin{minipage}[t]{.485\linewidth}
\textbf{Задание 2.}\quad$ f(x)=a^{x+b} $\\[1mm]
\begin{tikzpicture}
\begin{axis}[
 width=\linewidth,height=5.2cm,scale only axis=false,
 xmin=-0.78,xmax=2.6,ymin=0,ymax=26,
 xtick={0,1,2},ytick={0,5,10,15,20,25},axis lines=middle,
 grid=both,grid style={gridcolor,very thin},
 tick label style={font=\scriptsize},
 xlabel={$x$},ylabel={$y$},
 label style={font=\scriptsize},
 xlabel style={at={(axis description cs:1.01,.19)},anchor=west},
 ylabel style={at={(axis description cs:.40,1)},anchor=south},
 clip=true]
\addplot[accent,very thick,domain=-0.78:2.6,samples=200]{pow(3,2-x)};
\addplot[only marks,mark=*,mark size=2.4pt,accent] coordinates {(2,1) (0,9)};
\end{axis}
\end{tikzpicture}\\[-1mm]
\footnotesize Точки $A(2;1)$ и $B(0;9)$.\\
Найдите $x$, если $f(x)=27$.
\end{minipage}\par\vspace{12mm}
\textcolor{muted}{После решения проверьте обе точки подстановкой в восстановленную формулу. Только затем сверьтесь с ответами.}
\clearpage
\section*{08. Ответы и контрольные коэффициенты}
\sectionline
\textcolor{muted}{Проверяйте ответы после выполнения всех заданий. Значения $a$ и $b$ даны для диагностики ошибок.}
\begin{center}
\begin{tabularx}{\linewidth}{@{}c c c >{\raggedright\arraybackslash}X@{}}
\toprule
\textbf{№}&\textbf{$a$}&\textbf{$b$}&\textbf{Ответ}\\
\midrule
1&$4$&$-2$&$f(3)=62$\\
2&$\dfrac13$&$-2$&$x=-1$\\
\bottomrule
\end{tabularx}
\end{center}
\vspace{9mm}
\subsection*{Если ответ не совпал}
\begin{itemize}
\item Ошибка в $b$: заново подставьте точку с $x=0$ или $y=1$.
\item Ошибка в $a$: проверьте чётную и отрицательную степень.
\item Верные коэффициенты, но неверный ответ: повторите финальную подстановку и правила степеней.
\end{itemize}
\vfill
\begin{center}
\textcolor{darkaccent}{\textbf{Лёвин Артём Александрович}}\\[1mm]
\textcolor{muted}{эксклюзивно для Анны Банновой}
\end{center}
\end{document}
